%% NIH Biosketch — John Ehrlinger
%% Uses nihbiosketch.cls by Paul Magwene
%% https://github.com/pmagwene/latex-nihbiosketch
%%
%% Compile with XeLaTeX:
%%   xelatex JohnEhrlinger-Biosketch-NIH.tex
%%
%% nihbiosketch.cls must be in the same directory or on your TeX path.
%% Download: curl -fsSL https://raw.githubusercontent.com/pmagwene/latex-nihbiosketch/master/nihbiosketch.cls -o nihbiosketch.cls

%% nihbiosketch.cls calls \setmainfont{Arial} during class loading, which
%% causes a fatal fontspec error on systems without Arial (e.g. Linux CI).
%% quiet mode demotes that fatal error to a warning so the class loads;
%% the preamble below then sets the real font.
\PassOptionsToPackage{quiet}{fontspec}
\documentclass{nihbiosketch}

%% Use Arial when available (Mac), fall back to TeX Gyre Heros on CI
%% (metrically compatible, OpenType, ships with TeX Live Full).
\IfFontExistsTF{Arial}{%
  \setmainfont[Ligatures=TeX]{Arial}%
}{%
  \setmainfont[Ligatures=TeX]{TeX Gyre Heros}%
}

%% Personal information — auto-inserted as header via \piinfo
\name{Ehrlinger, John}
\eracommons{EHRLINJ}
\position{Assistant Staff -- Lead Data Scientist}
%% ORCID — add \orcid{} if supported by your version of nihbiosketch.cls,
%% otherwise include manually in the header or cover page.
%% ORCID iD: 0000-0002-5340-5154  https://orcid.org/0000-0002-5340-5154

\begin{document}

%% ================================================================
%% Education
%% ================================================================
\begin{education}
Case Western Reserve University, Cleveland, OH  & BS  & 05/1993 & Mechanical Engineering \\
Case Western Reserve University, Cleveland, OH  & MS  & 05/1994 & Mechanical / Aerospace Engineering \\
Case Western Reserve University, Cleveland, OH  & PhD & 05/2011 & Statistics \\
\end{education}

%% ================================================================
\section*{A. Personal Statement}
%% ================================================================

\begin{statement}
I am an applied statistician and data scientist with expertise in computational
statistical machine learning, deep learning, and their application to clinical
and engineering data. My research career spans both the theoretical development
of machine learning methods and their direct application to health outcomes,
with a particular focus on time-to-event data, longitudinal and temporal data
analysis, and survival methods.

My foundational methodological work on regularization and gradient
boosting---specifically the theoretical characterization of
$\ell_2$Boosting and its connections to the lasso and stagewise
regression---established a principled framework for understanding regularized
machine learning in high-dimensional settings. I have extended this work to
multivariate longitudinal settings through boosted tree ensembles, enabling
richer modeling of complex clinical trajectories. In parallel, I have developed
open-source R software tools (ggRandomForests, hazard) that make these methods
accessible to the broader statistical and biomedical research community.

I have since returned to Cleveland Clinic's Heart, Vascular \& Thoracic
Institute as Assistant Staff -- Lead Data Scientist in the Department of
Cardiothoracic Surgery, where I lead a team of data engineers and data
scientists focused on cardiovascular outcomes research. In this role I
drive statistical methods research as applied to observational clinical
research, establish best practices in software engineering and reproducible
research, and implement process improvements to optimize the departmental
research pipeline from data collection through publication.

\bigskip
\noindent\textbf{Selected citations relevant to this application:}

\begin{enumerate}

\item \textbf{Ehrlinger J}, Ishwaran H. Characterizing $\ell_2$Boosting.
\textit{The Annals of Statistics.} \textbf{40}(2): 1074--1101. 2012.

\item Pande A, Li L, Rajeswaran J, \textbf{Ehrlinger J}, Blackstone EH,
Ishwaran H, Lauer MS. Boosted multivariate trees for longitudinal data.
\textit{Machine Learning.} \textbf{106}(2): 277--305. 2017.

\item Hurst T, Xanthopoulos A, \textbf{Ehrlinger J}, et al. Dynamic
prediction of LVAD pump thrombosis based on lactate dehydrogenase trends.
\textit{ESC Heart Failure.} \textbf{6}(5): 1005--1014. 2019.

\item \textbf{Ehrlinger J}. ggRandomForests: Exploring Random Forest
Survival. \textit{arXiv preprint} arXiv:1612.08974 [stat.CO]. 2016.

\end{enumerate}
\end{statement}

%% ================================================================
\section*{B. Positions, Scientific Appointments, and Honors}
%% ================================================================

\subsection*{Positions and Scientific Appointments}

\begin{datetbl}
1994--1995    & Mechanical Engineer, Life Systems, Inc., Beachwood, OH \\
1995--1998    & Aerodynamic Engineer, Cooper Turbocompressor, Inc., Buffalo, NY \\
1999--2012    & Lead Systems Analyst: Scientific Programmer, Dept.\ of Thoracic and
                Cardiovascular Surgery, Cleveland Clinic, Cleveland, OH \\
2012--2015    & Assistant Staff, Dept.\ of Quantitative Health Sciences, Lerner Research
                Institute, Cleveland Clinic, Cleveland, OH \\
2015          & Assistant Professor of Medicine, Cleveland Clinic Lerner College of
                Medicine -- Case Western Reserve University \\
2015--2023    & Senior Data and Applied Scientist, Artificial Intelligence, Azure Global
                Commercial Industry (AGCI-AI), Microsoft, Cambridge, MA \\
2023--2024    & Senior Data Scientist -- Technical Lead, Altamira Technologies,
                McLean, VA \\
2024--present & Assistant Staff -- Lead Data Scientist, Dept.\ of Cardiothoracic Surgery,
                Heart, Vascular \& Thoracic Institute, Cleveland Clinic, Cleveland, OH \\
\end{datetbl}

\subsection*{Honors}

\begin{datetbl}
1992--1994  & NASA Graduate Research Fellowship \\
2003        & Cleveland Clinic Innovations Award \\
\end{datetbl}

%% ================================================================
\section*{C. Contributions to Science}
%% ================================================================

\begin{enumerate}

\item \textbf{Theoretical Foundations of Regularized Statistical Machine Learning.}

My dissertation and early research established rigorous theoretical properties
of gradient boosting under an $\ell_2$ loss function, characterizing the
functional form and convergence behavior of $\ell_2$Boosting in terms of
well-known shrinkage estimators including the lasso and LARS. I extended this
work to multivariate longitudinal settings via boosted tree ensembles, enabling
flexible nonparametric modeling of clinical time-course data.

\begin{enumerate}
\item \textbf{Ehrlinger J}, Ishwaran H. Characterizing $\ell_2$Boosting.
\textit{The Annals of Statistics.} \textbf{40}(2): 1074--1101. 2012.

\item Pande A, Li L, Rajeswaran J, \textbf{Ehrlinger J}, Blackstone EH,
Ishwaran H, Lauer MS. Boosted multivariate trees for longitudinal data.
\textit{Machine Learning.} \textbf{106}(2): 277--305. 2017.
\end{enumerate}

\item \textbf{Random Forest Methods and Open-Source Statistical Software.}

I have contributed to the development and dissemination of random forest
methods for survival, regression, and classification settings via the
ggRandomForests R package, and maintain the hazard R package, an open-source
implementation of multi-phase hazard analysis methods for time-to-event
decomposition currently available only through a SAS-dependent implementation.

\begin{enumerate}
\item \textbf{Ehrlinger J}. ggRandomForests: Visually Exploring a Random
Forest for Regression. \textit{arXiv preprint} arXiv:1501.07196 [stat.CO]. 2015.

\item \textbf{Ehrlinger J}. ggRandomForests: Exploring Random Forest
Survival. \textit{arXiv preprint} arXiv:1612.08974 [stat.CO]. 2016.
\end{enumerate}

\item \textbf{Temporal and Longitudinal Modeling for Cardiac Outcomes.}

A core application of my methodological work has been the modeling of
time-varying clinical data from cardiac patients, including parametric nonlinear
temporal decomposition models and mixture approaches to characterize atrial
fibrillation recurrence and prevalence.

\begin{enumerate}
\item Rajeswaran J, Blackstone EH, \textbf{Ehrlinger J}, Thuita L, et al.
Probability of atrial fibrillation after ablation: Using a parametric nonlinear
temporal decomposition mixed effects model. \textit{Statistical Methods in
Medical Research.} \textbf{27}(1): 126--141. 2018.

\item Li L, Mao H, Ishwaran H, \textbf{Ehrlinger J}, et al. Estimating the
prevalence of atrial fibrillation from a three class mixture model for repeated
diagnoses. \textit{Biometrical Journal.} \textbf{59}(2): 331--343. 2017.

\item Blackstone EH, Chang HL, Rajeswaran J, \textbf{Ehrlinger J}, et al.
Biatrial maze procedure versus pulmonary vein isolation for atrial fibrillation
during mitral valve surgery. \textit{The Journal of Thoracic and Cardiovascular
Surgery.} \textbf{157}(1): 234--243.e9. 2019.

\item Hurst T, Xanthopoulos A, \textbf{Ehrlinger J}, et al. Dynamic
prediction of LVAD pump thrombosis based on lactate dehydrogenase trends.
\textit{ESC Heart Failure.} \textbf{6}(5): 1005--1014. 2019.
\end{enumerate}

\item \textbf{Machine Learning Applications in Mechanical Circulatory Support.}

My clinical collaboration work at Cleveland Clinic's Heart and Vascular
Institute yielded several high-impact contributions identifying risk factors and
predictive biomarkers for LVAD pump thrombosis. The 2014 NEJM paper documenting
an abrupt increase in LVAD thrombosis rates had immediate regulatory and
clinical implications nationally.

\begin{enumerate}
\item Starling RC, Moazami N, Silvestry SC, Ewald G, Rogers JG, Milano CA,
Rame JE, Acker MA, Blackstone EH, \textbf{Ehrlinger J}, et al. Unexpected
Abrupt Increase in Left Ventricular Assist Device Thrombosis. \textit{New
England Journal of Medicine.} \textbf{370}(1): 33--40. 2014.

\item Smedira NG, Blackstone EH, \textbf{Ehrlinger J}, Thuita L, Pierce CD,
Moazami N, Starling RC. Current risks of HeartMate II pump thrombosis.
\textit{The Journal of Heart and Lung Transplantation.} \textbf{34}(12):
1527--1534. 2015.

\item Wojnarski CM, Roselli EE, Idrees JJ, et al. Machine-learning phenotypic
classification of bicuspid aortopathy. \textit{The Journal of Thoracic and
Cardiovascular Surgery.} \textbf{155}(2): 461--469.e4. 2018.
\end{enumerate}

\item \textbf{Cardiovascular Surgical Outcomes and Decision Support.}

Spanning my career at Cleveland Clinic, I have contributed analytical expertise
to outcomes studies covering LVAD support, transcatheter aortic valve
replacement (TAVR/PARTNER trial), and esophageal cancer surgery.

\begin{enumerate}
\item Szeto WY, Svensson LG, Rajeswaran J, \textbf{Ehrlinger J}, et al.
Appropriate Patient Selection or Healthcare Rationing? Lessons from Surgical
Aortic Valve Replacement in the PARTNER-I Trial. \textit{The Journal of
Thoracic and Cardiovascular Surgery.} \textbf{150}(3): 557--568. 2015.

\item Raja S, Rice TW, \textbf{Ehrlinger J}, et al. Importance of Residual
Cancer after Induction Therapy for Esophageal Adenocarcinoma. \textit{The
Journal of Thoracic and Cardiovascular Surgery.} \textbf{152}(3): 756--761.
2016.

\item Dalton JE, Glance LG, Mascha EJ, \textbf{Ehrlinger J}, Chamoun N,
Sessler DI. Impact of Present-on-admission Indicators on Risk-adjusted Hospital
Mortality Measurement. \textit{Anesthesiology.} \textbf{118}(6): 1298--1306.
2013.

\item Wojnarski CM, Svensson LG, Roselli EE, et al. Aortic Dissection in
Patients with Bicuspid Aortic Valve--Associated Aneurysm. \textit{Annals of
Thoracic Surgery.} \textbf{100}(5): 1666--1674. 2015.
\end{enumerate}

\end{enumerate}

%% ================================================================
\section*{D. Additional Information: Research Support and/or Scholastic
Performance}
%% ================================================================

\subsection*{Pending Research Support}

%% \grantinfo{source}{dates}{effort}{title}{PI}{role}
\grantinfo{NIH -- R01 (submitted Feb 2026)}
          {2026--}
          {N/A}
          {Improving Waitlist and Transplant Survival in Advanced Heart Failure
           Populations}
          {Multiple PIs: Eileen Hsich, M.D.\ and Eugene H. Blackstone, M.D.}
          {Co-Investigator}

\subsection*{Completed Research Support}

\grantinfo{NIH -- 1R01HL103552-01A1}
          {8/2011--5/2014}
          {N/A}
          {Ancillary Comparative Effectiveness of Atrial Fibrillation Ablation
           Surgery}
          {P.I.: Eugene H. Blackstone, M.D.}
          {Co-Investigator}

\grantinfo{State of Ohio -- Wright Center of Innovation}
          {8/2005--12/2010}
          {N/A}
          {Atrial Fibrillation Innovation Center State of Ohio}
          {P.I.: A. Marc Gillinov, M.D.}
          {Scientific programming and computing systems management}

\grantinfo{NIH -- HHSN26820080026C}
          {9/2008--9/2010}
          {N/A}
          {NHLBI Quantitative Clinical Cardiovascular Epidemiology Projects}
          {P.I.: Eugene H. Blackstone, M.D.}
          {Scientific programming and computing systems management}

\grantinfo{NIH -- 5R01HL072771}
          {10/2005--9/2008}
          {N/A}
          {Logical Analysis of Data and Cardiac Surgery Risk}
          {P.I.: Michael S.\ Lauer, M.D.}
          {Scientific programming and computing systems management}

\end{document}
